\chapter{Dimensi Hingga}\label{ch:b1:findim}

Sebuah ruang yang direntang oleh berhingga banyak vektor membawa
\emph{dimensi} yang terdefinisi dengan baik --- yaitu ukuran bersama semua basisnya. Adapun bukti
di bawah semuanya mengalir dari satu mesin kombinatorik, yaitu lema pertukarannya:
bahwa \omterm{def:b1:vspaces:free}{keluarga bebas} tak pernah dapat melebihi cacah keluarga pembangun. Bersama dimensinya
datanglah perkakas yang dipakai di mana-mana sesudahnya: yaitu teorema
pelengkapan \omterm{def:b1:vspaces:free}{basis}, \omterm{def:b1:findim:rank}{rank} sebuah keluarga, dan rumus Grassmann.

\section{Keberadaan basis}

\begin{definition}\label{def:b1:findim:def}
Ruang $E$ disebut \emph{berdimensi hingga}\index{dimensi hingga} bila ada padanya
keluarga pembangun yang hingga. (Kalau tidak ia \emph{berdimensi tak hingga}:
demikianlah $K[X]$, yang keluarga hingganya hanya merentang \omterm{def:b1:poly:def}{polinomial} yang derajatnya
terbatas.)
\end{definition}

\begin{theorem}[Lema pertukaran]\label{thm:b1:findim:exchange}
Misalkan $(g_1, \dots, g_n)$ membangun $E$ dan $(f_1, \dots, f_p)$ \omterm{def:b1:vspaces:free}{bebas}
di $E$. Maka $p \leq n$.
\end{theorem}

\begin{proof}
Kita membuktikannya lewat induksi pada $k \leq p$: bahwa \emph{setelah menomori ulang
$g_j$-nya, keluarga $(f_1, \dots, f_k, g_{k+1}, \dots, g_n)$ membangun
$E$} --- yang memaksa $k \leq n$ pada setiap tahapnya, sehingga $p \leq n$ pada
akhirnya.

Untuk $k = 0$: inilah hipotesisnya. Adapun langkahnya: anggaplah ia berlaku bagi $k - 1$; maka $f_k$
merupakan kombinasi $(f_1, \dots, f_{k-1},$ $g_k, \dots, g_n)$. Pada
kombinasi ini suatu $g_j$ (dengan $j \geq k$) berkoefisien taknol
---
sebab kalau tidak $f_k$ akan menjadi kombinasi $f_1, \dots, f_{k-1}$,
yang bertentangan dengan kebebasannya. Lalu nomorilah ulang sehingga $j = k$, lalu pecahkanlah bagi
$g_k$: maka $g_k$ merupakan kombinasi $(f_1, \dots, f_k, g_{k+1}, \dots,
g_n)$. Sehingga setiap vektor $E$, yang dinyatakan lewat keluarga beranggota $(k-1)$-nya,
lalu dapat dinyatakan ulang lewat keluarga beranggota $k$-nya: jadi ia membangun.
(Adapun jika $k - 1 = n$, tak ada $g$ yang tersisa sehingga $f_k$ akan menjadi kombinasi
$f_i$-nya belaka: yang mustahil; jadi $k \leq n$.)
\end{proof}

\begin{example}[Pertukarannya, yang diamati sekali]
\label{ex:b1:findim:exchangerun}
Di dalam $\R^2$, ambillah keluarga pembangun $\bigl(g_1, g_2\bigr) =
\bigl((1,0), (0,1)\bigr)$ dan \omterm{def:b1:vspaces:free}{keluarga bebas} $\bigl(f_1,
f_2\bigr) = \bigl((1,2), (3,4)\bigr)$. Langkah $1$: $f_1 = 1\cdot
g_1 + 2\cdot g_2$; koefisien $g_2$-nya taknol, sehingga
tukarkanlah $g_2$ dengan $f_1$: maka keluarga $\bigl(f_1, g_1\bigr)$
tetap membangun (karena $g_2 = \frac12(f_1 - g_1)$). Langkah $2$: $f_2 =
(3,4) = 2\,f_1 + 1\cdot g_1$; koefisien $g_1$ yang tersisa
bernilai taknol (dan memang harus: karena $f_2$ bukan kelipatan $f_1$),
sehingga tukarkanlah lagi: maka $\bigl(f_1, f_2\bigr)$ membangun $\R^2$. Seandainya
ada vektor \omterm{def:b1:vspaces:free}{bebas} ketiga $f_3$, tak ada $g$ yang tersisa untuk
menyerapnya --- dan persis begitulah lemanya melarang $3$ vektor
\omterm{def:b1:vspaces:free}{bebas} di $\R^2$. Jadi bukti di atas merupakan pembukuan ini yang dikerjakan
secara umum.
\end{example}

\begin{theorem}[Basis dalam dimensi hingga]\label{thm:b1:findim:bases}
Misalkan $E \neq \{0\}$ \omterm{def:b1:findim:def}{berdimensi hingga}.
\begin{enumerate}
  \item Dari setiap keluarga pembangun yang hingga kita dapat menyarikan sebuah \omterm{def:b1:vspaces:free}{basis}.
  \item (Teorema pelengkapan \omterm{def:b1:vspaces:free}{basis}\index{teorema pelengkapan basis})
        Setiap \omterm{def:b1:vspaces:free}{keluarga bebas} terperluas menjadi sebuah \omterm{def:b1:vspaces:free}{basis}, dengan memakai vektor sebarang
        keluarga pembangun yang dipilih.
  \item Semua \omterm{def:b1:vspaces:free}{basis} $E$ bersifat hingga, dengan cacah unsur yang
        sama: yaitu \emph{dimensinya} $\dim E$. (Kesepakatan:
        $\dim\{0\} = 0$.)
\end{enumerate}
\end{theorem}

\begin{proof}
(1) Buanglah, satu demi satu, sebarang vektor yang merupakan kombinasi
vektor lainnya; maka keluarganya tetap membangun, karena pada sebarang ungkapan
yang memakai vektor yang dibuang itu kita boleh menyulihkan kombinasinya atas
yang bertahan. Prosesnya berakhir --- karena setiap langkahnya mengecilkan sebuah
keluarga hingga sebanyak satu --- dan ia berhenti persis ketika tak ada vektor
yang tersisa yang merupakan kombinasi vektor lainnya. Adapun keluarga akhirnya tetap
membangun, dan ia \omterm{def:b1:vspaces:free}{bebas}: karena sebuah kombinasi nol yang taksepele akan
membawa suatu koefisien taknol, lalu \omterm{def:b1:arith:divides}{membaginya} dengan itu akan memecahkannya
bagi vektor yang bersesuaian dalam vektor lainnya, sehingga membuatnya
terbuang juga --- yang bertentangan dengan bahwa prosesnya sudah
berhenti.

(2) Misalkan $(f_1, \dots, f_p)$ \omterm{def:b1:vspaces:free}{bebas}, dan $(g_1, \dots, g_n)$ membangun.
Jalanilah $g_1, \dots, g_n$, sambil menambahkan $g_j$ pada keluarga yang berjalan
setiap kali ia belum berada di \omterm{def:b1:vspaces:span}{rentangnya}
(\cref{prop:b1:vspaces:freecriteria}~(2) menjaga keluarganya tetap \omterm{def:b1:vspaces:free}{bebas}).
Maka keluarga akhirnya \omterm{def:b1:vspaces:free}{bebas}, dan membangun: karena setiap $g_j$ terletak di
\omterm{def:b1:vspaces:span}{rentangnya} --- entah ia ditambahkan, entah ia sudah menjadi sebuah kombinasi.

(3) Dua \omterm{def:b1:vspaces:free}{basis} masing-masing \omterm{def:b1:vspaces:free}{bebas} dan masing-masing membangun: sehingga lema pertukarannya
memberikan kedua ketaksamaan antara kardinalitasnya. (Adapun kehinggaannya:
sebuah \omterm{def:b1:vspaces:free}{basis} bersifat \omterm{def:b1:vspaces:free}{bebas}, sehingga menurut pertukarannya ia tak lebih besar daripada keluarga
pembangun yang hingga.)
\end{proof}

\begin{example}\label{ex:b1:findim:dims}
$\dim K^n = n$ (dengan \omterm{def:b1:vspaces:free}{basis} kanoniknya); $\dim K_n[X] = n + 1$ (dengan monomialnya);
$\dim_\R \C = 2$; dan ruang penyelesaian $y'' + ay' + by = 0$ berdimensi
$2$ (\cref{thm:b1:diffeq:homogeneous2}: karena penyelesaiannya
terparameterkan secara \omterm{def:b1:logic:inj}{bijektif} dan linear oleh $(\lambda, \mu) \in K^2$).
\end{example}

\begin{example}[Himpunan yang sama, dua dimensi]
\label{ex:b1:findim:twodimensions}
\omterm{def:b1:logic:sets}{Himpunan} $\C^2$ berisi pasangan bilangan kompleks merupakan \omterm{def:b1:vspaces:def}{ruang vektor} atas $\C$
berdimensi $2$ (dengan \omterm{def:b1:vspaces:free}{basis} kanoniknya $e_1, e_2$) --- sekaligus
\omterm{def:b1:vspaces:def}{ruang vektor} atas $\R$ berdimensi $4$, dengan basisnya
\[
(1, 0),\quad (\iu, 0),\quad (0, 1),\quad (0, \iu):
\]
karena setiap $(z, w) = (a + \iu b,\ c + \iu d)$ mempunyai \omterm{prop:b1:vspaces:coordinates}{koordinat} real
$(a, b, c, d)$, secara tunggal. Jadi dimensi bukan sifat \omterm{def:b1:logic:sets}{himpunan}
vektornya semata: karena ia mencacah derajat kebebasan \emph{relatif
terhadap skalar yang diizinkan}, dan memaruhkan pasokan skalarnya dari $\C$
ke $\R$ melipatduakan cacahnya. (Adapun soal akhir pekannya memanfaatkan
kasus ekstrem kepekaan ini, dengan skalarnya disusutkan sampai
ke $\Q$.)
\end{example}

\begin{example}[Menjalankan algoritma pelengkapannya]
\label{ex:b1:findim:completion}
Lengkapilah \omterm{def:b1:vspaces:free}{keluarga bebas} $\bigl((1,1,1)\bigr)$ menjadi sebuah \omterm{def:b1:vspaces:free}{basis}
$\R^3$ dengan memakai vektor kanoniknya. Jalankanlah bukti
\cref{thm:b1:findim:bases}~(2) pada keluarga pembangun $(e_1,
e_2, e_3)$: apakah $e_1 \in \operatorname{Vect}(1,1,1)$? Bukan (karena kelipatan
$(1,1,1)$ berkoordinat sama) --- jadi tambahkanlah ia. Apakah $e_2 \in
\operatorname{Vect}\bigl((1,1,1), e_1\bigr)$? Sebuah kombinasi
$\alpha(1,1,1) + \beta e_1$ berkoordinat kedua dan ketiga yang sama,
sedangkan $e_2$ tidak --- jadi tambahkanlah ia. Maka keluarga
$\bigl((1,1,1), e_1, e_2\bigr)$ \omterm{def:b1:vspaces:free}{bebas} dengan $3$ vektor: berhentilah,
karena ia sebuah \omterm{def:b1:vspaces:free}{basis} (\cref{prop:b1:findim:twoofthree} akan meresmikan
refleks ini). Perhatikanlah bahwa hasilnya bergantung pada urutan
pemindaian $g_j$-nya: jadi pelengkapannya sebuah algoritma, bukan rumus.
\end{example}

\begin{proposition}[Aturan dua dari tiga]\label{prop:b1:findim:twoofthree}
Misalkan $\dim E = n$ dan $\mathcal{F}$ sebuah keluarga berisi tepat $n$ vektor
$E$. Maka
\[
\mathcal{F} \text{ bebas} \iff \mathcal{F} \text{ membangun} \iff
\mathcal{F} \text{ basis}.
\]
Lebih lanjut sebarang \omterm{def:b1:vspaces:free}{keluarga bebas} beranggota $\leq n$ vektor, dan sebarang keluarga pembangun
$\geq n$.
\end{proposition}

\begin{proof}
Adapun batas kardinalitasnya adalah lema pertukaran terhadap sebuah \omterm{def:b1:vspaces:free}{basis}. Jika
$\mathcal{F}$ (berukuran $n$) \omterm{def:b1:vspaces:free}{bebas} tetapi tak membangun, maka suatu $x$ terletak
di luar \omterm{def:b1:vspaces:span}{rentangnya}; sehingga menambahkan $x$ memberikan \omterm{def:b1:vspaces:free}{keluarga bebas} beranggota $n + 1$
vektor: yang mustahil. Jika $\mathcal{F}$ membangun tetapi tak \omterm{def:b1:vspaces:free}{bebas},
maka menyarikan sebuah \omterm{def:b1:vspaces:free}{basis} (\cref{thm:b1:findim:bases}~(1)) memberikan \omterm{def:b1:vspaces:free}{basis} beranggota
$< n$ vektor: yang mustahil.
\end{proof}

\begin{example}[Dua dari tiga, yang menghemat separuh kerjanya]
\label{ex:b1:findim:twoofthreerun}
Apakah $\bigl((1,1,0), (0,1,1), (1,0,1)\bigr)$ merupakan \omterm{def:b1:vspaces:free}{basis} $\R^3$?
Cacahlah: tiga vektor, dimensi tiga --- sehingga kebebasannya saja yang
memutuskan. Adapun sebuah kombinasi nol memberikan $a + c = 0$, $a + b = 0$, $b +
c = 0$; lalu menjumlahkan ketiganya, $2(a + b + c) = 0$, dan mengurangkan
setiap persamaan aslinya dari $a + b + c = 0$ menyisakan $b = c = a =
0$: jadi \omterm{def:b1:vspaces:free}{bebas}, sehingga sebuah \omterm{def:b1:vspaces:free}{basis}, dengan paruh pembangun
pemeriksaannya dipasok cuma-cuma oleh teoremanya. Bandingkanlah
\cref{ex:b1:vspaces:testbasis}, yang di situ pemeriksaan ganda yang sama
harus dikerjakan dengan tangan --- jadi satu bab teori
tertukar persis menjadi penghematan itu, pada setiap pemeriksaan \omterm{def:b1:vspaces:free}{basis} bagi
sisa buku ini.
\end{example}

\begin{method}[Menghitung sebuah dimensi]
\label{met:b1:findim:computedim}
Ada tiga rute baku, dengan urutan kekerapan yang menurun.
\begin{enumerate}
  \item \emph{Parameterkanlah, lalu bacalah sebuah \omterm{def:b1:vspaces:free}{basis}.} Pecahkanlah
        kendala pendefinisinya, nyatakanlah unsur umumnya
        secara linear dalam parameter yang bertahan, lalu periksalah bahwa
        vektor yang mengalikan parameternya bersifat \omterm{def:b1:vspaces:free}{bebas}: maka
        dimensinya adalah cacah parameternya. (Dijalankan di bawah pada
        sebuah \omterm{def:b1:vspaces:subspace}{subruang} $\R^4$ yang konkret.)
  \item \emph{Tunjukkanlah sebuah parameterisasi linear yang \omterm{def:b1:logic:inj}{bijektif}.} Ketika
        unsurnya ditentukan oleh berhingga banyak nilai ---
        yaitu syarat awal sebuah rekurensi
        (\cref{exo:b1:findim:10}), atau koefisien sebuah rumus
        penyelesaian (\cref{ex:b1:findim:dims}) --- maka dimensinya adalah
        cacah nilai itu.
  \item \emph{Pakailah rumusnya.} Yaitu Grassmann bagi irisan dan
        jumlah, \omterm{def:b1:findim:rank}{rank} bagi \omterm{def:b1:vspaces:span}{rentang}, dan nanti rank--nulitas bagi
        kernel dan peta: karena dimensi biasanya
        \emph{dihitung}, bukan ditebak.
\end{enumerate}
Pada ketiga rutenya, aturan dua dari tiganyalah \omterm{def:b1:topology:closure}{penutupnya}:
begitu cacahnya cocok, kebebasan atau pembangunnya saja sudah menyimpulkan.
\end{method}

\begin{example}[Rute 1, selengkapnya]
\label{ex:b1:findim:parametrize}
Dimensi $H = \{(x, y, z, t) \in \R^4 : x + y + z + t = 0
\text{ dan } x = t\}$. Pecahkanlah: $t = x$ dan $y + z = -2x$, sehingga $z =
-2x - y$ dengan $x, y$ yang \omterm{def:b1:vspaces:free}{bebas}:
\[
(x,\ y,\ -2x - y,\ x) = x\,(1, 0, -2, 1) + y\,(0, 1, -1, 0) .
\]
Kedua vektornya \omterm{def:b1:vspaces:free}{bebas} (lihatlah dua \omterm{prop:b1:vspaces:coordinates}{koordinat} pertamanya:
$(x, y) = (0,0)$), sehingga keduanya membentuk \omterm{def:b1:vspaces:free}{basis} $H$ dan $\dim H = 2$.
Cacahnya memang terramalkan --- karena dua kendala linear yang saling \omterm{def:b1:vspaces:free}{bebas}
di $\R^4$ semestinya masing-masing memakan satu dimensi --- tetapi
parameterisasinyalah yang membuktikannya \emph{dan} menyerahkan sebuah \omterm{def:b1:vspaces:free}{basis}, yang
tak pernah dilakukan ramalannya semata; adapun \cref{exo:b1:findim:9} mengubah
semboyan ``setiap persamaan memakan paling banyak satu dimensi'' menjadi sebuah
teorema.
\end{example}

\begin{example}[Rute 2, selengkapnya]
\label{ex:b1:findim:route2}
Dimensi $W = \{P \in \R_n[X] : P(1) = P(2) = 0\}$ (untuk $n
\geq 2$). Menurut teorema faktor yang diterapkan dua kali
(\cref{thm:b1:poly:factor}; karena akar $1$ dan $2$-nya berbeda),
$P \in W$ tepat ketika $P = (X - 1)(X - 2)\,Q$ dengan $\deg Q \leq
n - 2$. Adapun korespondensi $Q \mapsto (X-1)(X-2)Q$ bersifat linear,
mencapai seluruh $W$, dan \omterm{def:b1:logic:inj}{injektif} (karena hasil kalinya nol hanya bila
$Q = 0$): sehingga $W$ terparameterkan secara \omterm{def:b1:logic:inj}{bijektif} dan linear oleh
$\R_{n-2}[X]$, jadi
\[
\dim W = \dim \R_{n-2}[X] = n - 1 .
\]
Adapun sebuah \omterm{def:b1:vspaces:free}{basis} datang bersama parameterisasinya: yaitu peta
monomialnya, $\bigl((X-1)(X-2),\ (X-1)(X-2)X,\ \dots,\
(X-1)(X-2)X^{n-2}\bigr)$. Jadi setiap kendala penilaian baru pada sebuah
titik yang segar berongkos tepat satu dimensi --- yaitu tulang punggung
pencacahan \omterm{thm:b1:poly:lagrange}{interpolasi Lagrange}
(\cref{thm:b1:poly:lagrange}).
\end{example}

\begin{example}[Sebuah kendala integral pun berongkos satu dimensi]
\label{ex:b1:findim:integralconstraint}
Dimensi $H = \{P \in \R_2[X] : \int_0^1 P = 0\}$. Dengan menulis
$P = a + bX + cX^2$, kendalanya berbunyi $a + \frac b2 + \frac
c3 = 0$; lalu pecahkanlah bagi $a$ dan parameterkanlah:
\[
P = b\Bigl(X - \frac12\Bigr) + c\Bigl(X^2 - \frac13\Bigr),
\]
sehingga $H = \operatorname{Vect}\bigl(X - \frac12,\ X^2 -
\frac13\bigr)$, berdimensi $2$ (karena kedua \omterm{def:b1:poly:def}{polinomialnya} berderajat
berbeda: jadi \omterm{def:b1:vspaces:free}{bebas}). Jadi satu syarat linear --- entah sebuah penilaian,
sebuah \omterm{thm:b1:integration:def}{integral}, atau sebarang resep linear lainnya --- membuang paling banyak satu
dimensi, dan tepat satu begitu syaratnya tak
identik nol. Adapun \cref{ch:b1:linmaps} akan menamai resep semacam itu
\emph{bentuk linear} dan \omterm{def:b1:logic:sets}{himpunan} penyelesaiannya \emph{hiperbidang};
sedangkan vektor \omterm{def:b1:vspaces:free}{basis} yang ditemukan di sini muncul kembali pada
\cref{ch:b1:euclid} sebagai awal keluarga Legendre.
\end{example}

\section{Subruang, rank, Grassmann}

\begin{theorem}[Subruang]\label{thm:b1:findim:subspaces}
Misalkan $E$ \omterm{def:b1:findim:def}{berdimensi hingga} dan $F$ sebuah \omterm{def:b1:vspaces:subspace}{subruang}. Maka $F$
\omterm{def:b1:findim:def}{berdimensi hingga}, $\dim F \leq \dim E$, dengan kesamaannya jika dan hanya
jika $F = E$. Lebih lanjut setiap \omterm{def:b1:vspaces:subspace}{subruang} mempunyai \omterm{def:b1:vspaces:subspace}{subruang} \omterm{def:b1:vspaces:sum}{pelengkap}.
\end{theorem}

\begin{proof}
\omterm{def:b1:vspaces:free}{Keluarga bebas} $F$ beranggota paling banyak $\dim E$ vektor (menurut lema pertukaran
di $E$: karena \omterm{def:b1:vspaces:free}{keluarga bebas} $F$ khususnya \omterm{def:b1:vspaces:free}{bebas} di $E$, dan $E$
mempunyai keluarga pembangun yang hingga). Di antara \omterm{def:b1:vspaces:free}{keluarga bebas} $F$
pilihlah satu yang berukuran \emph{maksimal} $p$ --- yang mungkin karena ukurannya
berupa bilangan bulat yang terbatas oleh $\dim E$. Maka ia membangun $F$: sebab kalau tidak suatu
$x \in F$ akan terletak di luar \omterm{def:b1:vspaces:span}{rentangnya}, dan menambahkan $x$ akan memberikan
\omterm{def:b1:vspaces:free}{keluarga bebas} $F$ berukuran $p + 1$
(\cref{prop:b1:vspaces:freecriteria}~(2)), yang bertentangan dengan
kemaksimalannya. Jadi karena \omterm{def:b1:vspaces:free}{bebas} dan membangun, ia sebuah \omterm{def:b1:vspaces:free}{basis} $F$, dan
$\dim F = p \leq \dim E$. Adapun jika $p = \dim E = n$: \omterm{def:b1:vspaces:free}{keluarga bebas} beranggota $n$ vektor
$E$ merupakan \omterm{def:b1:vspaces:free}{basis} $E$ (\cref{prop:b1:findim:twoofthree}), sehingga $F
\supseteq \operatorname{span} = E$. Untuk \omterm{def:b1:vspaces:sum}{pelengkapnya}: lengkapilah sebuah \omterm{def:b1:vspaces:free}{basis}
$(f_1, \dots, f_p)$ bagi $F$ menjadi \omterm{def:b1:vspaces:free}{basis} $(f_1, \dots, f_p, g_{p+1},
\dots, g_n)$ bagi $E$ (menurut teorema pelengkapan \omterm{def:b1:vspaces:free}{basis}); maka $G =
\operatorname{Vect}(g_{p+1}, \dots, g_n)$ memenuhi $E = F \oplus G$
(karena keberadaan dan ketunggalan penguraiannya = \omterm{prop:b1:vspaces:coordinates}{koordinatnya} pada
\omterm{def:b1:vspaces:free}{basis} besarnya).
\end{proof}

\begin{definition}[Rank sebuah keluarga]\label{def:b1:findim:rank}
Adapun \emph{rank}\index{rank!keluarga} sebuah keluarga hingga berisi
vektor adalah dimensi \omterm{def:b1:vspaces:span}{rentangnya}:
$\operatorname{rk}(x_1, \dots, x_p) = \dim
\operatorname{Vect}(x_1, \dots, x_p) \leq \min(p, \dim E)$,
dengan kesamaannya dengan $p$ jika dan hanya jika keluarganya \omterm{def:b1:vspaces:free}{bebas}.
\end{definition}

\begin{example}[Menghitung rank lewat penghapusan]
\label{ex:b1:findim:rankcomputation}
\omterm{def:b1:findim:rank}{Rank} $\bigl((1,2,3), (2,3,4), (3,4,5), (1,1,1)\bigr)$ di
$\R^3$. Adapun \omterm{def:b1:vspaces:span}{rentangnya} tak berubah ketika kita mengurangkan dari sebuah vektor sebuah
kombinasi vektor lainnya (karena kedua keluarganya merentang kombinasi yang
sama): jadi gantilah $(2,3,4)$ dengan $(2,3,4) - (1,2,3) = (1,1,1)$
dan $(3,4,5)$ dengan $(3,4,5) - (1,2,3) = (2,2,2)$. Maka \omterm{def:b1:vspaces:span}{rentangnya} kini
$\operatorname{Vect}\bigl((1,2,3), (1,1,1), (2,2,2),
(1,1,1)\bigr) = \operatorname{Vect}\bigl((1,2,3), (1,1,1)\bigr)$,
dan kedua vektornya tak sebanding: sehingga \omterm{def:b1:findim:rank}{ranknya} $2$. Adapun
prosedur kurangi-dan-buang ini disistematiskan sebagai penghapusan
Gauss pada \cref{ch:b1:det}.
\end{example}

\begin{example}[Menjumlahkan lewat perangkaian]
\label{ex:b1:findim:sumconcat}
Ambillah, di dalam $\R^3$,
\[
F = \operatorname{Vect}\bigl((1,2,3),\ (1,1,1)\bigr),
\qquad
G = \operatorname{Vect}\bigl((2,3,4)\bigr) .
\]
Adapun jumlahnya $F + G$ direntang oleh keluarga terangkai atas ketiga
pembangunnya, dan
\[
(2, 3, 4) = (1, 2, 3) + (1, 1, 1)
\]
menunjukkan bahwa yang ketiganya berlebihan: sehingga $F + G = F$, berdimensi $2$ ---
setara dengan itu $G \subseteq F$, yang dipertunjukkan relasinya.
Adapun Grassmann menegaskannya: $\dim(F \cap G) = 2 + 1 - 2 = 1 = \dim G$.
Jadi jumlahnya dihitung dengan merangkai pembangunnya lalu mengurangi
tumpukannya lewat algoritma \omterm{def:b1:findim:rank}{ranknya}; sehingga tak pernah diperlukan teknik yang baru.
\end{example}

\begin{theorem}[Rumus Grassmann]\label{thm:b1:findim:grassmann}
Untuk \omterm{def:b1:vspaces:subspace}{subruang} $F, G$ pada $E$ yang \omterm{def:b1:findim:def}{berdimensi
hingga}:\index{rumus Grassmann}
\[
\dim(F + G) = \dim F + \dim G - \dim(F \cap G) .
\]
Khususnya $F + G$ bersifat langsung jika dan hanya jika $\dim(F + G) = \dim F + \dim G$.
\end{theorem}

\begin{proof}
Mulailah dari sebuah \omterm{def:b1:vspaces:free}{basis} $(e_1, \dots, e_r)$ bagi $F \cap G$; lalu lengkapilah ia
menjadi \omterm{def:b1:vspaces:free}{basis} $(e_1, \dots, e_r, f_1, \dots, f_s)$ bagi $F$ dan menjadi
\omterm{def:b1:vspaces:free}{basis} $(e_1, \dots, e_r, g_1, \dots, g_t)$ bagi $G$
(\cref{thm:b1:findim:bases}~(2)). Kita mengklaim bahwa
\[
\mathcal{B} = (e_1, \dots, e_r, f_1, \dots, f_s, g_1, \dots, g_t)
\]
merupakan \omterm{def:b1:vspaces:free}{basis} $F + G$; sehingga rumusnya menyusul lewat pencacahan: $(r + s) +
(r + t) - r = r + s + t$.

Adapun $\mathcal{B}$ membangun $F + G$: karena sebarang $u + v$ (dengan $u \in F$, $v \in G$)
terjabarkan lewatnya. Sedangkan kebebasannya: andaikan $\sum \alpha_i e_i + \sum
\beta_j f_j + \sum \gamma_k g_k = 0$. Maka vektor $w = \sum \gamma_k
g_k = -\sum\alpha_i e_i - \sum\beta_j f_j$ terletak di $G \cap F$, sehingga ia
terjabarkan pada $(e_i)$ belaka; padahal $w$ juga terjabarkan pada $(g_k)$, dan pada
\omterm{def:b1:vspaces:free}{basis} $G$ kedua ungkapan itu harus berimpit: sehingga semua
$\gamma_k = 0$ (dan \omterm{prop:b1:vspaces:coordinates}{koordinat} $e$-nya cocok). Lalu relasinya tereduksi
menjadi $\sum\alpha_i e_i + \sum\beta_j f_j = 0$, yaitu relasi pada \omterm{def:b1:vspaces:free}{basis}
$F$: sehingga semua koefisien yang tersisa lenyap.
\end{proof}

\begin{example}\label{ex:b1:findim:planes}
Dua bidang $F, G$ yang berbeda (berdimensi $2$) pada $\R^3$ memenuhi $F + G
= \R^3$ (karena jumlahnya tegas memuat sebuah bidang), sehingga $\dim(F \cap G) =
2 + 2 - 3 = 1$: jadi keduanya selalu berpotongan sepanjang sebuah garis --- sehingga tak ada
``bidang yang sejajar'' lewat titik asalnya.
\end{example}

\begin{example}[Grassmann dalam kerja, di $\R^4$]
\label{ex:b1:findim:grassmannrun}
Misalkan $F = \operatorname{Vect}\bigl(e_1,\ e_2,\ (1,1,1,0)\bigr)$
dan $G = \operatorname{Vect}(e_3, e_4)$ di $\R^4$. Adapun dimensinya:
$\dim F = 3$ (karena pembangun ketiganya berkoordinat ketiga taknol,
di luar $\operatorname{Vect}(e_1, e_2)$) dan $\dim G = 2$. Adapun jumlahnya:
$F + G$ memuat $e_1, e_2, e_4$ dan $e_3 = (1,1,1,0) - e_1 -
e_2$: sehingga ia seluruh $\R^4$. Lalu Grassmann \emph{menghitung}
ukuran irisannya tanpa penghapusan sama sekali:
\[
\dim(F \cap G) = 3 + 2 - 4 = 1 .
\]
Untuk mengenali garisnya, lihatlah di dalam $G$: sebuah vektor $(0, 0, c, d)$
terletak di $F$ tepat ketika ia $a e_1 + b e_2 + \lambda(1,1,1,0)$,
yang memaksa $\lambda = -a = -b$ dan $d = 0$: sehingga irisannya adalah
$\operatorname{Vect}(e_3)$ --- yang selaras, karena $e_3$ tadi
ditunjukkan di $F$ dan terletak di $G$ menurut definisinya. Jadi khaslah
pembagian kerjanya: Grassmann meramalkan \emph{seberapa banyak} yang harus dicari,
lalu sistem linearnya menemukan \emph{apanya}.
\end{example}

\begin{example}[Mengiriskan lewat persamaan]
\label{ex:b1:findim:intersectequations}
Ketika kedua \omterm{def:b1:vspaces:subspace}{subruangnya} datang sebagai \omterm{def:b1:logic:sets}{himpunan} penyelesaian, mengiriskannya sekadar
menumpuk persamaannya. Di dalam $\R^3$: $F = \{x + y + z = 0\}$ dan
$G = \{x = y\}$ memberikan
\[
F \cap G = \{x = y,\ 2x + z = 0\}
= \{(x,\ x,\ -2x)\} = \operatorname{Vect}(1, 1, -2),
\]
yaitu sebuah garis. Adapun periksa silangnya lewat Grassmann: $F + G = \R^3$ (karena bidangnya
berbeda, sehingga jumlahnya tegas memuat sebuah bidang), jadi
$\dim(F\cap G) = 2 + 2 - 3 = 1$. Jadi kedua pemerian sebuah
\omterm{def:b1:vspaces:subspace}{subruang} --- lewat persamaan, dan lewat pembangun --- masing-masing membuat satu
operasinya sepele: karena persamaannya beririsan lewat penumpukan, sedangkan pembangunnya
berjumlah lewat perangkaian; dan mengalihkan antarkeduanya persis merupakan apa yang
dimaksud memecahkan sistem linear (\cref{ch:b1:det}).
\end{example}

\begin{remark}[Jebakan yang lazim]
\label{rem:b1:findim:pitfalls}
\emph{Dimensi tak terjumlahkan sepanjang jumlahnya} kecuali jika jumlahnya langsung:
karena dua bidang $\R^3$ mempunyai $\dim(F + G) = 3$, bukan $4$; jadi koreksilah selalu
lewat suku irisannya (menurut Grassmann).
\emph{Sebuah inklusi menuntut dimensinya \emph{dan} inklusinya}:
karena $\dim F = \dim G$ semata tak pernah memberikan $F = G$ (dengan dua garis berbeda
pada $\R^2$); jadi kasus kesamaan pada
\cref{thm:b1:findim:subspaces} menuntut $F \subseteq G$ lebih dulu.
\emph{Mencacah parameter belumlah sebuah bukti}: karena ``dua persamaan di
$\R^4$, jadi berdimensi $2$'' gagal ketika persamaannya
bergantungan (sebab $x + y = 0$ dan $2x + 2y = 0$ menyisakan dimensi $3$);
sehingga hanya parameterisasi yang sungguhan atau perhitungan \omterm{def:b1:findim:rank}{rank} yang memutuskan.
\emph{Jangan berbicara tentang dimensi sesuatu yang bukan \omterm{def:b1:vspaces:subspace}{subruang}}: karena \omterm{def:b1:logic:sets}{himpunan}
penyelesaian sistem yang \emph{takhomogen} kehilangan $0$; sehingga ``dimensinya''
adalah dimensi ruang penyelesaian homogen yang bersesuaian
(\cref{ch:b1:det} memerincinya).
\emph{Dimensi tak hingga itu ada}: karena $K[X]$ memuat \omterm{def:b1:vspaces:free}{keluarga bebas}
berukuran apa pun (yaitu monomialnya), sehingga tak ada keluarga pembangun yang hingga
yang mungkin --- jadi \omterm{def:b1:logic:statement}{pernyataan} seperti aturan dua dari tiganya
tegas \omterm{def:b1:findim:def}{berdimensi hingga} dan gagal dengan buruk pada $K[X]$
(\cref{cor:b1:linmaps:samedim} akan menunjukkan hal yang sama bagi
keinjektifan/kesurjektifannya).
\emph{\omterm{def:b1:findim:rank}{Rank} menyangkut \omterm{def:b1:vspaces:span}{rentangnya}, bukan daftarnya}: karena mengulang sebuah vektor,
mengurutkannya ulang, atau menskalakannya dengan konstanta taknol membiarkan \omterm{def:b1:findim:rank}{ranknya}
tak berubah, dan $\operatorname{rk} = p$ (yaitu cacah vektornya)
merupakan \emph{sifat yang harus dibuktikan} --- yang persis merupakan kebebasannya. Adapun
sebuah keluarga berisi $5$ vektor dengan \omterm{def:b1:findim:rank}{rank} $2$ membawa kelebihan senilai tiga
vektor, yang dilokalkan penghapusannya
(\cref{ex:b1:findim:rankcomputation}) secara eksplisit.
\end{remark}

\begin{remark}[Ke mana dimensinya bekerja]
\label{rem:b1:findim:whereused}
Dimensi merupakan argumen pencacahan kesayangan buku ini mulai dari sini.
\cref{ch:b1:linmaps} membuktikan teorema rank--nulitas, yaitu
versi fungsional rumus Grassmann; \cref{ch:b1:matrices}
menghitung \omterm{def:b1:findim:rank}{ranknya} lewat reduksi baris; sedangkan \cref{ch:b1:det} mengubah ``$n$
vektor $K^n$ membentuk \omterm{def:b1:vspaces:free}{basis}'' menjadi satu bilangan yang taknol. Adapun
soal akhir pekan di bawah memperlihatkan dimensi mengerjakan \emph{aritmetika}:
karena mencacah dimensi di atas \omterm{def:b1:structures:field}{lapangan} $\Q$ membuktikan \omterm{def:b1:logic:statement}{pernyataan}
keirasionalan yang tampak tak tersentuh dengan tangan. Pada jilid Tahun~3
pencacahan dimensi yang sama, yang dihaluskan oleh teori \omterm{def:b1:structures:group}{grup}, memutuskan
masalah konstruksi klasik mana yang terpecahkan --- dan kisah itu
adalah teori Galois.
\end{remark}

\begin{remark}[Cakrawala di dalam Buku 3]
\label{rem:b1:findim:perspectives}
Dimensi merupakan besaran kekal sisa jilid ini,
dan hukum kekekalannya pantas dinamai lebih dulu.
\cref{ch:b1:linmaps} membuktikan $\dim E = \dim\ker u +
\operatorname{rk} u$: bahwa apa yang digilas sebuah \omterm{def:b1:logic:map}{pemetaan} linear ditambah apa yang
disimpannya selalu berjumlah sumbernya. \cref{ch:b1:det} menghaluskan ini
menjadi struktur \omterm{def:b1:logic:sets}{himpunan} penyelesaiannya: bahwa $p$ anu dikurangi
$\operatorname{rk} A$ poros menyisakan dimensi
ruang penyelesaiannya, yang ditunjukkan penghapusan Gauss sebagai parameter
yang \omterm{def:b1:vspaces:free}{bebas}. \cref{ch:b1:euclid} membelah $\dim E = \dim F + \dim
F^\perp$ secara ortogonal, sedangkan soal akhir pekan
\cref{ch:b1:multivar} membelanjakan persis anggaran itu: yaitu $n$ titik
data, $2$ parameter yang dicocokkan, dan $n - 2$ dimensi sisanya.
Jadi setiap kali sebuah cacah menolak berimbang pada bab berikutnya,
galatnya berupa kernel yang terlupakan atau jumlah yang tak langsung --- jadi kembalilah ke
rumus Grassmann lebih dulu.
\end{remark}

\section{Latihan}

\begin{exercise}[$\star$]\label{exo:b1:findim:1}
Berikanlah sebuah \omterm{def:b1:vspaces:free}{basis} beserta dimensi dari:
\begin{enumerate}
  \item $F = \{(x,y,z) \in \R^3 : x + y + z = 0\}$;
  \item $G = \{(x,y,z,t) \in \R^4 : x = y,\ z = 2t\}$;
  \item $H = \{P \in \R_3[X] : P(1) = P'(1) = 0\}$.
\end{enumerate}
\end{exercise}

\begin{exercise}[$\star$]\label{exo:b1:findim:2}
Hitunglah \omterm{def:b1:findim:rank}{rank} keluarga
$\bigl((1,1,1), (1,2,3), (3,5,7), (0,1,2)\bigr)$ di $\R^3$, lalu
sarikanlah sebuah \omterm{def:b1:vspaces:free}{basis} bagi \omterm{def:b1:vspaces:span}{rentangnya}.
\end{exercise}

\begin{exercise}[$\star$]\label{exo:b1:findim:3}
Lengkapilah \omterm{def:b1:vspaces:free}{keluarga bebas} $\bigl((1,1,0,0), (0,0,1,1)\bigr)$ menjadi
\omterm{def:b1:vspaces:free}{basis} $\R^4$ dengan memakai vektor kanoniknya, lalu benarkanlah.
\end{exercise}

\begin{exercise}[$\star$]\label{exo:b1:findim:4}
Buktikan bahwa $\bigl(1, X, X(X-1), X(X-1)(X-2)\bigr)$ merupakan \omterm{def:b1:vspaces:free}{basis}
$\R_3[X]$, lalu carilah \omterm{prop:b1:vspaces:coordinates}{koordinat} $X^3$ padanya.
\end{exercise}

\begin{exercise}[$\star\star$]\label{exo:b1:findim:5}
Misalkan $F$ dan $G$ \omterm{def:b1:vspaces:subspace}{subruang} berdimensi $4$ dan $5$ pada sebuah ruang
$E$ dengan $\dim E = 7$. Apa saja nilai yang mungkin bagi $\dim(F \cap
G)$? Berikanlah sebuah kejadian yang mewujudkan setiap nilainya dengan $E = \R^7$.
\end{exercise}

\begin{exercise}[$\star\star$]\label{exo:b1:findim:6}
Misalkan $H = \{P \in \R_n[X] : P(1) = 0\}$. Buktikan bahwa $H$ merupakan
hiperbidang $\R_n[X]$ (yaitu \omterm{def:b1:vspaces:subspace}{subruang} berdimensi $n$), tunjukkanlah sebuah
\omterm{def:b1:vspaces:free}{basis} $H$ \emph{(pikirkanlah teorema faktornya: $P = (X-1)Q$)}, lalu
berikanlah sebuah garis \omterm{def:b1:vspaces:sum}{pelengkapnya}.
\end{exercise}

\begin{exercise}[$\star\star$]\label{exo:b1:findim:7}
Misalkan $u_1, \dots, u_p$ vektor dengan \omterm{def:b1:findim:rank}{rank} $r$. Buktikan bahwa membuang
satu vektornya menghasilkan keluarga dengan \omterm{def:b1:findim:rank}{rank} $r$ atau $r - 1$, dan bahwa
menambahkan satu vektor menghasilkan \omterm{def:b1:findim:rank}{rank} $r$ atau $r + 1$. Lalu simpulkan bahwa \omterm{def:b1:findim:rank}{ranknya}
berubah paling banyak $1$ di bawah sebarang penyisipan atau penghapusan tunggal.
\end{exercise}

\begin{exercise}[$\star\star\star$]\label{exo:b1:findim:8}
Misalkan $F_1 \subseteq F_2 \subseteq \dots \subseteq F_k$ \omterm{def:b1:vspaces:subspace}{subruang}
$E$ (dengan $\dim E = n$) dan $F_i \neq F_{i+1}$ untuk setiap $i$. Buktikan $k
\leq n + 1$. Lalu simpulkan bahwa rantai \omterm{def:b1:vspaces:subspace}{subruang} $\R^n$ yang naik tegas
berpanjang paling banyak $n + 1$, lalu tunjukkanlah satu yang berpanjang
maksimal.
\end{exercise}

\begin{exercise}[$\star\star\star$]\label{exo:b1:findim:9}
Misalkan $E$ berdimensi $n$ dan $F$, $G$ dua hiperbidang (berdimensi
$n - 1$), dengan $F \neq G$. Hitunglah $\dim(F \cap G)$. Samaratakanlah: bahwa
irisan $k$ hiperbidang berdimensi $\geq n - k$.
\end{exercise}

\begin{exercise}[$\star\star$]\label{exo:b1:findim:10}
Misalkan $E$ \omterm{def:b1:logic:sets}{himpunan} barisan real yang memenuhi $u_{n+2} =
3u_{n+1} - 2u_n$ untuk setiap $n$.
\begin{enumerate}
  \item Tunjukkan bahwa $E$ merupakan \omterm{def:b1:vspaces:subspace}{subruang} ruang barisannya, dan
        bahwa sebuah barisan $E$ sepenuhnya ditentukan, secara linear,
        oleh pasangan $(u_0, u_1)$; lalu simpulkanlah $\dim E = 2$.
  \item Periksalah bahwa barisan konstan $(1)$ dan barisan geometri
        $(2^n)$ terletak di $E$ dan membentuk \omterm{def:b1:vspaces:free}{basis} $E$.
  \item Carilah barisan $E$ dengan $u_0 = 0$, $u_1 = 1$.
\end{enumerate}
\end{exercise}

\begin{exercise}[$\star\star$]\label{exo:b1:findim:11}
Misalkan $E$ berdimensi $n$.
\begin{enumerate}
  \item Jika $F, G$ \omterm{def:b1:vspaces:subspace}{subruang} dengan $\dim F + \dim G > n$, buktikanlah
        $F \cap G \neq \{0\}$. Perlihatkanlah: bahwa dua \omterm{def:b1:vspaces:subspace}{subruang}
        berdimensi $51$ dan $50$ pada $\R^{100}$ selalu berbagi sebuah
        vektor taknol.
  \item Jika $H$ sebuah hiperbidang dan $F$ \omterm{def:b1:vspaces:subspace}{subruang} dengan $F \cap H =
        \{0\}$, buktikanlah $\dim F \leq 1$.
\end{enumerate}
\end{exercise}

\begin{exercise}[$\star\star\star$]\label{exo:b1:findim:12}
(\omterm{def:b1:vspaces:sum}{Pelengkap} bersama) Misalkan $F, G$ \omterm{def:b1:vspaces:subspace}{subruang} $E$ (yang \omterm{def:b1:findim:def}{berdimensi
hingga}) dengan $\dim F = \dim G$. Buktikan bahwa $F$ dan $G$ mempunyai
\omterm{def:b1:vspaces:sum}{pelengkap} yang \emph{bersama}: yakni ada \omterm{def:b1:vspaces:subspace}{subruang} $S$ dengan $E = F
\oplus S = G \oplus S$. \emph{(Berinduksilah menurun pada $\dim F$: jika $F
\neq G \neq E$, pilihlah $x \notin F \cup G$ ---
\cref{exo:b1:vspaces:12} mengizinkannya --- lalu tinjaulah $F \oplus
\operatorname{Vect}(x)$ dan $G \oplus \operatorname{Vect}(x)$.)}
\end{exercise}

\section{Soal: Hukum menara Dedekind}

\begin{problem}\label{pb:b1:findim:1}
Tak ada apa pun pada Bab 18--19 yang memakai sesuatu tentang skalarnya di luar
aksioma \omterm{def:b1:structures:field}{lapangannya} (\cref{def:b1:structures:field}): sehingga kita boleh
mengambil $K = \Q$ lalu mengukur \omterm{def:b1:logic:sets}{himpunan} \emph{bilangan real}
dengan tolok ukur dimensi atas $\Q$. Soal ini menghitung
dimensi $\Q(\sqrt2, \sqrt3)$, membuktikan \emph{\omterm{pb:b1:findim:1}{hukum menara}
Dedekind} $\dim_\Q M = \dim_\Q K \cdot \dim_K M$, lalu memanen
teorema keirasionalan lewat pencacahan dimensi belaka --- tanpa
$\varepsilon$, tanpa desimal, hanya dengan \omterm{def:b1:vspaces:free}{basis}.\index{hukum menara}
\index{lapangan!sublapangan $\R$}\index{bilangan irasional}

\medskip
\noindent\textbf{Bagian I --- Skalar yang rasional.}
\begin{enumerate}
  \item Periksalah bahwa $\Q$ sebuah \omterm{def:b1:structures:field}{lapangan} dan bahwa setiap definisi dan
        bukti pada Bab 18--19 hanya memakai aksioma \omterm{def:b1:structures:field}{lapangan} bagi
        skalarnya; lalu simpulkan bahwa $\R$ merupakan \omterm{def:b1:vspaces:def}{ruang vektor} atas $\Q$ dan
        bahwa lema pertukaran, teorema basisnya, dan
        rumus Grassmann berlaku atas $\Q$. Tunjuklah satu langkah
        pada bukti \cref{thm:b1:findim:exchange} yang di situ
        pembagian oleh skalar taknol dikerjakan.
  \item Tunjukkan bahwa $(1, \sqrt2)$ \omterm{def:b1:vspaces:free}{bebas} atas $\Q$ tetapi terikat atas
        $\R$. Tetapkanlah $\Q(\sqrt2) = \operatorname{Vect}_\Q(1,
        \sqrt2) = \{a + b\sqrt2 : a, b \in \Q\}$; berapakah
        $\dim_\Q \Q(\sqrt2)$?
  \item Tunjukkan bahwa $\Q(\sqrt2)$ stabil terhadap perkalian.
        Untuk $u = a + b\sqrt2$ tetapkanlah $\sigma(u) = a - b\sqrt2$ dan
        $N(u) = u\,\sigma(u) = a^2 - 2b^2$. Tunjukkan $\sigma(uv) =
        \sigma(u)\sigma(v)$, lalu simpulkan $N(uv) = N(u)N(v)$, lalu tunjukkan
        $N(u) \neq 0$ setiap kali $u \neq 0$.
  \item Simpulkan bahwa setiap $u \in \Q(\sqrt2)$ yang taknol mempunyai
        inversnya di $\Q(\sqrt2)$, yakni $u^{-1} =
        \sigma(u)/N(u)$: sehingga $\Q(\sqrt2)$ merupakan sublapangan $\R$.
        Hitunglah $\dfrac1{3 + 2\sqrt2}$ dan $\dfrac1{1 + \sqrt2}$.
\end{enumerate}

\medskip
\noindent\textbf{Bagian II --- Menyertakan $\sqrt3$.}
\begin{enumerate}[resume]
  \item Buktikan bahwa $\sqrt6 \notin \Q$ \emph{(bandingkanlah eksponen
        $2$ pada kedua ruas $6q^2 = p^2$, seperti pada
        \cref{exo:b1:arith:7})}, lalu bahwa $\sqrt3 \notin
        \Q(\sqrt2)$ \emph{(kuadratkanlah $\sqrt3 = a + b\sqrt2$ lalu
        bahaslah kasus $ab \neq 0$, $b = 0$, $a = 0$)}.
  \item Misalkan $M = \operatorname{Vect}_\Q(1, \sqrt2, \sqrt3,
        \sqrt6)$. Tunjukkan bahwa $M$ stabil terhadap perkalian
        \emph{(sebuah tabel hasil kali vektor basisnya sudah
        mencukupi)}.
  \item Tulislah $K = \Q(\sqrt2)$. Tunjukkan bahwa $(1, \sqrt3)$ \omterm{def:b1:vspaces:free}{bebas}
        \emph{atas \omterm{def:b1:structures:field}{lapangan} $K$}, lalu simpulkan bahwa $M = K +
        K\sqrt3$ merupakan \omterm{def:b1:vspaces:def}{ruang vektor} atas $K$ berdimensi $2$ dengan
        \omterm{def:b1:vspaces:free}{basis} $(1, \sqrt3)$.
  \item Buktikan bahwa $(1, \sqrt2, \sqrt3, \sqrt6)$ \omterm{def:b1:vspaces:free}{bebas} atas
        $\Q$, sehingga $\dim_\Q M = 4$. \emph{(Kelompokkanlah sebuah relasi nol
        sebagai $(a + b\sqrt2) + (c + d\sqrt2)\sqrt3 = 0$ lalu terapkan
        pertanyaan 7 kemudian 2.)}
\end{enumerate}

\medskip
\noindent\textbf{Bagian III --- \omterm{pb:b1:findim:1}{Hukum menaranya}.} Misalkan $\Q \subseteq K
\subseteq M \subseteq \R$ dengan $K$ dan $M$ sublapangan, $(e_1,
\dots, e_m)$ sebuah \omterm{def:b1:vspaces:free}{basis} $K$ sebagai \omterm{def:b1:vspaces:def}{ruang vektor} atas $\Q$, dan $(f_1,
\dots, f_n)$ sebuah \omterm{def:b1:vspaces:free}{basis} $M$ sebagai \omterm{def:b1:vspaces:def}{ruang vektor} atas $K$.
\begin{enumerate}[resume]
  \item Tunjukkan bahwa $mn$ hasil kali $(e_i f_j)$ membangun $M$ atas
        $\Q$.
  \item Tunjukkan bahwa keluarga $(e_i f_j)$ \omterm{def:b1:vspaces:free}{bebas} atas $\Q$.
        \emph{(Tatalah ulang sebuah kombinasi nol atas $\Q$ sebagai $\sum_j
        \bigl(\sum_i \lambda_{ij} e_i\bigr) f_j$, yang koefisien
        dalamnya hidup di $K$.)}
  \item Tutuplah dengan \emph{\omterm{pb:b1:findim:1}{hukum menara} Dedekind}:
        \[
        \dim_\Q M \;=\; \dim_\Q K \,\cdot\, \dim_K M ,
        \]
        lalu periksalah ia pada $\Q \subseteq \Q(\sqrt2) \subseteq M$
        terhadap pertanyaan 7 dan 8.
  \item (\omterm{def:b1:structures:field}{Lapangan} secara cuma-cuma) Misalkan $A \subseteq \R$ sebuah
        \omterm{def:b1:vspaces:subspace}{subruang} atas $\Q$ yang \omterm{def:b1:findim:def}{berdimensi hingga} yang memuat $1$ dan
        stabil terhadap perkalian, dan misalkan $u \in A$, $u \neq
        0$. Tunjukkan bahwa jika $(a_1, \dots, a_d)$ sebuah \omterm{def:b1:vspaces:free}{basis} $A$,
        maka $(u a_1, \dots, u a_d)$ kembali menjadi \omterm{def:b1:vspaces:free}{basis} $A$;
        lalu simpulkan bahwa $u$ mempunyai invers \emph{di $A$}: sehingga $A$ merupakan
        sublapangan $\R$. Pertanyaan sebelumnya yang mana yang dipulihkan ini?
  \item Tunjukkan bahwa untuk setiap $x \in A$ (dengan, seperti pada pertanyaan 12, $\dim_\Q
        A = d$) keluarga $(1, x, x^2, \dots, x^{d})$ bersifat terikat:
        sehingga setiap unsur $A$ merupakan akar sebuah \omterm{def:b1:poly:def}{polinomial} taknol
        berkoefisien rasional, yang berderajat paling banyak $d$.
\end{enumerate}

\medskip
\noindent\textbf{Bagian IV --- Satu bilangan membangkitkan segalanya.}
Tetapkanlah $s = \sqrt2 + \sqrt3$.
\begin{enumerate}[resume]
  \item Hitunglah \omterm{prop:b1:vspaces:coordinates}{koordinat} $s^2$, $s^3$ dan $s^4$ pada
        \omterm{def:b1:vspaces:free}{basis} $(1, \sqrt2, \sqrt3, \sqrt6)$ bagi $M$.
  \item Tunjukkan bahwa $(1, s, s^2, s^3)$ \omterm{def:b1:vspaces:free}{bebas} atas $\Q$, lalu simpulkanlah
        $\operatorname{Vect}_\Q(1, s, s^2, s^3) = M$: bahwa setiap
        unsur $M$ merupakan \omterm{def:b1:poly:def}{polinomial} rasional dalam $s$. Nyatakanlah
        $\sqrt2$ dan $\sqrt3$ sebagai \omterm{def:b1:poly:def}{polinomial} yang demikian.
  \item Periksalah $s^4 - 10s^2 + 1 = 0$, lalu tunjukkanlah bahwa $X^4 - 10X^2
        + 1$ merupakan \emph{\omterm{def:b1:poly:def}{polinomial monik} berderajat terkecil}
        yang lenyap di $s$. Tentukanlah keempat akar realnya.
  \item Simpulkan: $1/s = 10s - s^3$; lalu kenalilah bilangan ini. Tunjukkanlah
        bahwa $a + b\sqrt2 + c\sqrt3 + d\sqrt6$ (dengan $a,b,c,d \in \Q$)
        bersifat rasional jika dan hanya jika $b = c = d = 0$; sehingga
        khususnya $\sqrt2 + \sqrt3 + \sqrt6$ irasional
        (yang memperkuat \cref{exo:b1:reals:7}).
\end{enumerate}

\medskip
\noindent\textbf{Bagian V --- Akar pangkat tiganya tetap di luar.} Tetapkanlah $t
= 2^{1/3}$.
\begin{enumerate}[resume]
  \item Tunjukkan bahwa $X^3 - 2$ tak berakar rasional \emph{(lewat
        kriteria akar rasional pada \cref{exo:b1:poly:5}, atau lewat
        pencacahan valuasinya)}; lalu simpulkanlah bahwa $(1, t)$ \omterm{def:b1:vspaces:free}{bebas} atas
        $\Q$.
  \item Tunjukkan bahwa $(1, t, t^2)$ \omterm{def:b1:vspaces:free}{bebas} atas $\Q$. \emph{(Jika suatu
        $P \in \Q[X]$ taknol berderajat $\leq 2$ membunuh $t$, ambillah
        satu yang berderajat terkecil lalu bagilah $X^3 - 2$ dengannya,
        \cref{thm:b1:poly:division}; lalu simpulkanlah bahwa $X^3 - 2$
        akan berakar rasional.)} Lalu simpulkanlah bahwa $A =
        \operatorname{Vect}_\Q(1, t, t^2)$ berdimensi $3$,
        stabil terhadap perkalian, dan merupakan sublapangan $\R$.
  \item Buktikan bahwa $t \notin M$: sebab kalau tidak $M$ akan menjadi \omterm{def:b1:vspaces:def}{ruang
        vektor} atas \omterm{def:b1:structures:field}{lapangan} $A$, dan \omterm{pb:b1:findim:1}{hukum menaranya} akan memaksa
        $3 \mid 4$. Jadi $2^{1/3}$ bukan kombinasi rasional
        atas $1, \sqrt2, \sqrt3, \sqrt6$.
  \item Simpulkan bahwa $t \notin \Q(\sqrt2)$, bahwa $t + \sqrt2$
        irasional, dan bahwa tak ada bilangan rasional $a, b$ yang memenuhi
        $2^{1/3} = a + b\sqrt3$.
\end{enumerate}

\medskip
\noindent\textbf{Bagian VI --- Semua sublapangannya, dan sintesisnya.}
\begin{enumerate}[resume]
  \item Buktikanlah \omterm{def:b1:arith:divides}{keterbagian} derajatnya: bahwa jika $\Q \subseteq A
        \subseteq B \subseteq \R$ merupakan sublapangan dengan $\dim_\Q B$
        yang hingga, maka $\dim_\Q A$ \omterm{def:b1:arith:divides}{membagi} $\dim_\Q B$. Apa saja
        dimensi yang mungkin bagi sublapangan $M$?
  \item Tentukanlah \emph{semua} sublapangan $M$ yang berdimensi $2$
        atas $\Q$. \emph{(Reduksikanlah ke pencarian $w = a + b\sqrt2
        + c\sqrt3 + d\sqrt6 \in M \setminus \Q$ dengan $w^2 \in
        \Q$: jabarkanlah $w^2$ lalu lenyapkanlah ketiga \omterm{prop:b1:vspaces:coordinates}{koordinat}
        irasionalnya.)}
  \item Nyatakanlah $\dfrac1{1 + \sqrt2 + \sqrt3}$ pada \omterm{def:b1:vspaces:free}{basis} $(1,
        \sqrt2, \sqrt3, \sqrt6)$ \emph{(kalikanlah dengan sekawan
        yang terpilih baik)}.
  \item Sintesis, dalam empat kalimat: mengapakah dimensi atas $\Q$ merupakan
        invarian \emph{aritmetika} sebuah \omterm{def:b1:logic:sets}{himpunan} bilangan real; apa yang
        dipaksakan kehinggaan dimensinya atas setiap unsurnya
        (pertanyaan 13); bagaimanakah aturan dua dari tiganya menghasilkan
        invers entah dari mana (pertanyaan 12); dan apa yang
        dilarang \omterm{pb:b1:findim:1}{hukum menaranya} (pertanyaan 20). Namailah teorema yang dibuktikan
        pada Bagian III.

\end{enumerate}
\end{problem}
